\begingroup
\renewcommand{\arraystretch}{1.45}
\Needspace{34\baselineskip}
\InterpolationArticle{15}
It follows from the preceding theorem that, if the values of a function $X$ of the form $F$ for $x=a$, $b$, $c$, $d$ etc. are respectively $A$, $B$, $C$, $D$ etc., any other similar function of the arc $x$ that satisfies the same values must be contained in the formula $X+PY$, where by $P$ we denote the product
\[
\sin\tfrac12(x-a)\sin\tfrac12(x-b)\sin\tfrac12(x-c)\sin\tfrac12(x-d)\quad\text{etc.}
\]
and by $Y$ an undetermined function of the arc $x$ of the form $F$ or of the form
\SourceLeaf{298}
\(G\), according as the number of values \(a,b,c,d\) etc., which we will denote by \(\mu\), is even or odd. We now observe:

I. In Articles 10, 11 and 12 we explained how to find a function \(X\) of the \(m^{\mathrm{th}}\) order (that is, not extending beyond \(\cos mx\) and \(\sin mx\)) that satisfies the \(\mu\) given values \(A,B,C,D\) etc., with \(m\) \(=\frac{1}{2}\mu-\frac{1}{2}\) or \(\frac{1}{2}\mu\), according as \(\mu\) is odd or even. The product \(P\) is evidently of the \({\frac{1}{2}\mu}^{\mathrm{th}}\) order. Thus, when \(\mu\) is odd, so that \(Y\) is of the \({\frac{1}{2}}^{\mathrm{th}}\) order or higher, \(PY\) will be at least of the \({\frac{1}{2}\mu+\frac{1}{2}}^{\mathrm{th}}\) order and therefore of higher order than \(X\). It follows that functions other than \(X\) satisfying the same given values can only be of higher order. When \(\mu\) is even, however, \(PY\) is of the \({\frac{1}{2}\mu}^{\mathrm{th}}\) order only when a fixed quantity is taken for \(Y\); hence in this case there are indeed infinitely many functions of similar form and the same order as \(X\) that satisfy the same values, but they are all contained in the form \(X+hP\), where \(h\) denotes a fixed quantity. The same conclusions follow directly from the method by which, in Articles 10 and 11, we derived the function \(X\).

\Needspace{23\baselineskip}
II. Conversely, if a function \(X'\) were known from elsewhere that satisfies all the given values \(A,B,C,D\) etc., but extends to a higher order than is needed, the function \(Y\) can be determined so that \(X'+PY\) is reduced to the \({\frac{1}{2}\mu-\frac{1}{2}}^{\mathrm{th}}\) or \({\frac{1}{2}\mu}^{\mathrm{th}}\) order (according as \(\mu\) is odd or even). Let the highest terms in \(X'\) be \(K\cos nx+L\sin nx\), and those in \(P\) be \(k\cos\frac{1}{2}\mu x+l\sin\frac{1}{2}\mu x\), and take as the highest terms in \(Y\)
\[
\rho=-\frac{2(Kk+Ll)}{kk+ll}\cos(n-\tfrac{1}{2}\mu)x
     +\frac{2(Kl-Lk)}{kk+ll}\sin(n-\tfrac{1}{2}\mu)x
\]
Calculation will then show that the coefficients of \(\cos nx\) and \(\sin nx\) in \(X'+PY\) become \(=0\), provided that \(n>\frac{1}{2}\mu\); in the same way, from the highest terms of the expanded function \(X'+P\rho\) (which will be of the \({n-1}^{\mathrm{th}}\) order), the next lower terms in \(Y\) (of the \({n-\frac{1}{2}\mu-1}^{\mathrm{th}}\) order) are determined. This operation must be repeated until \(X'+PY\) has been reduced to an order no higher than \(\frac{1}{2}\mu\), and thus to the \({\frac{1}{2}\mu-\frac{1}{2}}^{\mathrm{th}}\) order when \(\mu\) is odd, or to the \({\frac{1}{2}\mu}^{\mathrm{th}}\) order when \(\mu\) is even. It is readily seen that this reduction cannot be carried further.

\Needspace{42\baselineskip}
III. Thus in the first case (when \(\mu\) is odd) we will necessarily arrive in this way at the same function \(X\) that we would have found from the given values \(A,B,C,D\) etc. by the method set out above. In the second case, however, the method just described supplies a function which is indeed of the same order as the function obtained by the method of Article 12, or by that of Article 11 when the last term is set
\SourceLeaf{299}
\(=0\), but may differ from it. Let the former function be \(X''\) and the latter \(X'''\); thus \(X'''=X''+hP\), where \(h\) denotes a fixed quantity. Let us take the highest terms in \(X''\) to be \(K'\cos\frac{1}{2}\mu x+L'\sin\frac{1}{2}\mu x\), and those in \(X'''\) to be \(K''\cos\frac{1}{2}\mu x+L''\sin\frac{1}{2}\mu x\), so that \(K''=K'+hk\), \(L''=L'+hl\). It is now not difficult to show from Article 11 or Article 12 that, if the sum \(a+b+c+d+\text{etc.}\) is denoted by \(s\), then \(K''\cos\frac{1}{2}s+L''\sin\frac{1}{2}s=0\), from which follows
\[
\begin{aligned}
h&=-\frac{K'\cos\frac{1}{2}s+L'\sin\frac{1}{2}s}
              {k\cos\frac{1}{2}s+l\sin\frac{1}{2}s}\\[5pt]
K''&=\phantom{-}\frac{K'l-L'k}
              {k\cos\frac{1}{2}s+l\sin\frac{1}{2}s}\sin\tfrac{1}{2}s\\[5pt]
L''&=-\frac{K'l-L'k}
              {k\cos\frac{1}{2}s+l\sin\frac{1}{2}s}\cos\tfrac{1}{2}s
\end{aligned}
\]
It is no harder to show that
\[
k=\pm\frac{\cos\frac{1}{2}s}{2^{\mu-1}},\quad
l=\pm\frac{\sin\frac{1}{2}s}{2^{\mu-1}}
\]
where the upper sign applies when \(\frac{1}{2}\mu\) is even, and the lower when it is odd. Hence, in the values just given for \(K''\) and \(L''\), the denominator becomes \(=\pm\frac{1}{2^{\mu-1}}\), and thus
\[
\begin{aligned}
K''&=\phantom{-}(K'\sin\tfrac{1}{2}s-L'\cos\tfrac{1}{2}s)\sin\tfrac{1}{2}s\\[5pt]
L''&=-(K'\sin\tfrac{1}{2}s-L'\cos\tfrac{1}{2}s)\cos\tfrac{1}{2}s
\end{aligned}
\]

\Needspace{42\baselineskip}
\InterpolationArticle{16}

So far we have considered functions containing both cosines and sines. Very often, however, one or the other is absent, and it is appropriate to treat these kinds of functions separately. Both cases could indeed be reduced to the general case considered so far; it nevertheless seems even more appropriate to reduce this investigation to the problem of Article 3.

\textit{First}, if it is known that all the coefficients \(\beta',\beta'',\beta'''\) etc. in the function \(X\) (Article 10) are \(=0\), the number of unknowns is reduced to \(m+1\), and it therefore suffices to know, for that many distinct values of \(x\), the corresponding values of the function \(X\). Since in this case the value of the function \(X\) is unchanged when \(x\) is replaced by \(-x\), it is evident that not only values of \(x\) whose difference equals a circumference or a multiple of it, but also those whose sum equals \(0\) or a circumference or a multiple of it, must not be regarded as distinct. Or, in general, for two values of \(x\) to count as distinct, their cosines must be un\SourceLeaf{300}%
equal. With this understood, let the \(m+1\) given values of the function \(X\) be \(A, B, C, D\) etc., corresponding to the values of \(x\) given by \(a, b, c, d\) etc.; a general expression is sought for the value \(T\) taken by the function for \(x=0\).
Since, in our case, the function \(X\) can also be brought into the form
\[
\gamma+\gamma'\cos x+\gamma''(\cos x)^2+\gamma'''(\cos x)^3+\text{etc.}+\gamma^m(\cos x)^m
\]
Article 3 gives
\[
\begin{aligned}
T={}&\frac{(\cos t-\cos b)(\cos t-\cos c)(\cos t-\cos d)\mathinner{.\,.\,.}}{(\cos a-\cos b)(\cos a-\cos c)(\cos a-\cos d)\mathinner{.\,.\,.}}\,A\\[5pt]
&+\frac{(\cos t-\cos a)(\cos t-\cos c)(\cos t-\cos d)\mathinner{.\,.\,.}}{(\cos b-\cos a)(\cos b-\cos c)(\cos b-\cos d)\mathinner{.\,.\,.}}\,B\\[5pt]
&+\frac{(\cos t-\cos a)(\cos t-\cos b)(\cos t-\cos d)\mathinner{.\,.\,.}}{(\cos c-\cos a)(\cos c-\cos b)(\cos c-\cos d)\mathinner{.\,.\,.}}\,C\\[5pt]
&+\frac{(\cos t-\cos a)(\cos t-\cos b)(\cos t-\cos c)\mathinner{.\,.\,.}}{(\cos d-\cos a)(\cos d-\cos b)(\cos d-\cos c)\mathinner{.\,.\,.}}\,D\\[5pt]
&+\text{etc.}
\end{aligned}
\]

\Needspace{43\baselineskip}
\textit{Second}, if all the coefficients \(\alpha,\alpha',\alpha'',\alpha'''\) etc. vanish, the number of unknowns is \(m\), and it therefore suffices to know the values of the function \(X\) for the same number of distinct values of \(x\). Since the value of the function \(X\) for \(x=-t\) follows directly from its value for
\(x=+t\) (merely by reversing the sign), these values cannot be regarded as distinct here either. It follows readily that here, as in the preceding case, only those values of \(x\) whose cosines are unequal count as distinct. Moreover, since the value of the function \(X\) for \(x=0,180^0,360^0\)
etc., and in general for a value of \(x\) whose sine is \(=0\), is automatically \(=0\)
and is therefore already given by the nature of the problem, we suppose that such a value is not counted among the \(m\) given values. Now, since it is known that
\[
\frac{\sin nx}{\sin x}=2\cos(n-1)x+2\cos(n-3)x+2\cos(n-5)x+\mathinner{.\,.\,.}+2\cos x
\]
for an even value of \(n\), and
\[
=2\cos(n-1)x+2\cos(n-3)x+2\cos(n-5)x+\mathinner{.\,.}+2\cos 2x+1
\]
for an odd value of \(n\), it follows readily that \(\frac{X}{\sin x}\) can be brought into the form
\[
\delta+\delta'\cos x+\delta''(\cos x)^2+\delta'''(\cos x)^3+\text{etc.}+\delta^{m-1}(\cos x)^{m-1}
\]
\SourceLeaf{301}
and hence, since the values of the function \(\frac{X}{\sin x}\) for \(x=a,b,c,d\mathinner{.\,.\,.}\) are
\(=\frac{A}{\sin a},\frac{B}{\sin b},\frac{C}{\sin c},\frac{D}{\sin d}\mathinner{.\,.\,.\,.}\), the same Article 3 gives
\[
\begin{aligned}
T={}&\frac{\sin t(\cos t-\cos b)(\cos t-\cos c)(\cos t-\cos d)\mathinner{.\,.\,.}}{\sin a(\cos a-\cos b)(\cos a-\cos c)(\cos a-\cos d)\mathinner{.\,.\,.}}\,A\\[5pt]
&+\frac{\sin t(\cos t-\cos a)(\cos t-\cos c)(\cos t-\cos d)\mathinner{.\,.\,.}}{\sin b(\cos b-\cos a)(\cos b-\cos c)(\cos b-\cos d)\mathinner{.\,.\,.}}\,B\\[5pt]
&+\frac{\sin t(\cos t-\cos a)(\cos t-\cos b)(\cos t-\cos d)\mathinner{.\,.\,.}}{\sin c(\cos c-\cos a)(\cos c-\cos b)(\cos c-\cos d)\mathinner{.\,.\,.}}\,C\\[5pt]
&+\frac{\sin t(\cos t-\cos a)(\cos t-\cos b)(\cos t-\cos c)\mathinner{.\,.\,.}}{\sin d(\cos d-\cos a)(\cos d-\cos b)(\cos d-\cos c)\mathinner{.\,.\,.}}\,D\\[5pt]
&+\text{etc.}
\end{aligned}
\]
In both cases, moreover, the formula for \(T\) becomes identical with the function \(X\) on replacing
\(t\) by \(x\), since it holds without restriction for any value of \(t\).

\endgroup
